\subsection{情感亲密的日常维护}

性不是技术动作的堆叠，而是情感连接的一种形态——同一场性爱，有无情感在场，体验天差地别。

\begin{itemize}
	\item \textbf{日常积累}：性外之情是性内之欲的燃料——日常的拥抱、牵手、真诚的对话与共同承担，都是"亲密账户"的存款；长期冷战的关系里，性很难凭技术起死回生
	\item \textbf{在场与回应}：亲密中的"情感在场"——眼神、回应、随对方节奏调整——比任何技巧都更被伴侣报告为"好的性"的核心要素
	\item \textbf{冲突的性与性后的冲突}：吵架后的"和解性爱"可以修复关系，也可能成为回避沟通的模式；区分"我们想亲近"与"我们不想谈问题"——前者健康，后者透支
	\item \textbf{无性危机}：当亲密消失（无性婚姻状态持续），先修情感通道再修性通道——从每天 20 分钟无手机交流开始，比直接"安排性爱"更有效
\end{itemize}
